% 第6章 6.1：三点路径上的顶点函数与边函数。数据：f=(3,1,2)；
% (a) 记录方向 e1: 1->2, e2: 2->3，B^T f=(f1-f2,f2-f3)=(2,-1)；
% (b) 只把 e2 反向为 3->2，边差变为 (2,1)，平方 (4,1) 不变。
\begin{tikzpicture}[x=1mm,y=1mm,>=Stealth,font=\figurefont\footnotesize,text=FigureInk,
  vertex/.style={circle,draw=FigureStroke,fill=FigureBlueFill,line width=.8pt,minimum size=6.5mm,inner sep=0pt},
  arc/.style={->,draw=FigureStroke,line width=.8pt,shorten >=.6mm,shorten <=.6mm},
  val/.style={text=FigureBlue},
  ebox/.style={draw=FigureYellow,fill=FigureYellowFill,rounded corners=1.8pt,inner sep=2pt,minimum width=8mm},
  note/.style={font=\FigureNoteFont\fontsize{7.5}{9}\selectfont,text=FigureMuted}]
\foreach \dx/\flip/\tag in {0/0/(a),60/1/(b)}{
 \begin{scope}[xshift=\dx mm]
  \node[vertex] (a) at (4,14) {$1$};
  \node[vertex] (b) at (24,14) {$2$};
  \node[vertex] (c) at (44,14) {$3$};
  \node[val] at (4,22) {$f_1=3$};
  \node[val] at (24,22) {$f_2=1$};
  \node[val] at (44,22) {$f_3=2$};
  \draw[arc] (a)--(b) node[midway,above=.3mm,note]{$e_1$};
  \ifnum\flip=0
   \draw[arc] (b)--(c) node[midway,above=.3mm,note]{$e_2$};
   \node[ebox] at (14,5) {$+2$};
   \node[ebox] at (34,5) {$-1$};
   \node[note] at (24,-3) {平方：$4$，$1$};
  \else
   \draw[arc,draw=FigureRed] (c)--(b) node[midway,above=.3mm,note]{$e_2$};
   \node[ebox] at (14,5) {$+2$};
   \node[ebox,draw=FigureRed,fill=FigureRedFill] at (34,5) {$+1$};
   \node[note] at (24,-3) {平方：$4$，$1$};
  \fi
  \node[anchor=west] at (-2,30) {\tag};
 \end{scope}}
\node[anchor=west] at (4,30) {$e_1:1\to2$，$e_2:2\to3$};
\node[anchor=west] at (64,30) {$e_1:1\to2$，$e_2$改为$3\to2$};
\end{tikzpicture}
